\documentclass[12pt,a4paper,oneside]{book}
\usepackage[utf8]{inputenc}
\usepackage[T1]{fontenc}
\usepackage[english,french]{babel}
\usepackage{geometry}
\usepackage[table]{xcolor}

\geometry{top=2.5cm,bottom=2.5cm,left=2.5cm,right=2.5cm}
\definecolor{bleuIA}{RGB}{37, 99, 235}

\begin{document}

\begin{titlepage}
\begin{center}
\vspace*{3cm}
{\Huge\bfseries\color{bleuIA}E-GRAPHISME}\\[0.3cm]
{\LARGE by ELECTRON}\\[2cm]
{\huge ENTERPRISE SECURITY CENTER}\\[0.3cm]
{\huge Cybersécurité • Gouvernance}\\[1cm]
{\Large Volume 17}\\[3cm]
{\Large Version 2026}
\end{center}
\end{titlepage}

\chapter{Vision}
La sécurité doit être intégrée à chaque composant.

\chapter{Zero Trust}
Ne jamais faire confiance par défaut.

\chapter{Identity \& Access Management}
Gestion utilisateurs, entreprises, administrateurs.

\chapter{Authentification}
Mot de passe, 2FA, clés de sécurité.

\chapter{Chiffrement}
Données au repos et en transit.

\chapter{Protection des Bases de Données}
Sécurisation des accès et sauvegardes.

\chapter{Protection Ollama}
Contrôle des modèles autorisés.

\chapter{Protection des Agents IA}
Identité, permissions, journal d'activité.

\chapter{Pare-feu}
Filtrage du trafic.

\chapter{Surveillance}
Détection d'intrusion.

\chapter{Journal d'Audit}
Traçabilité des actions.

\chapter{AI Security Guardian}
Agent IA de sécurité.

\chapter{Policy Center}
Gestion des politiques.

\chapter{Threat Intelligence}
Analyse des menaces.

\chapter{Trust Center}
Portail de confiance.

\end{document}
